\section{Reward Design}\label{sec:reward}

\subsection{State of the art}\label{sec:sota-reward}
\todo{Draft with verified references: reward shaping (potential-based shaping preserves optimal policies,
arbitrary shaping does not); reward hacking and specification gaming; safe RL (constrained MDPs, Lagrangian methods,
shielding). Consequence for this work: a large collision penalty only changes the \emph{expected return}; it is not a
guarantee of collision avoidance.}

\subsection{Reward configurations}\label{sec:reward-configs}

\paragraph{Design goals.}
The assignment asks for the default \env{HighwayEnv} reward and at least one alternative that rewards driving at an
appropriate speed, penalizes collisions, and penalizes undesirable lane positioning and unnecessary lane changes, with
three configurations: efficiency-oriented, safety-oriented and balanced. We therefore keep the environment's default
reward as a reference and define three custom variants that share one functional form and differ only in their
weights. All variants are used with both algorithms, and the weights were fixed before any agent was trained: they
were chosen from simple, interpretable considerations and checked only on scripted baseline policies, never tuned on
the learning curves of \ddqn{} or \ppo{}. This avoids the trial-and-error reward tuning that Knox et
al.~\cite{knox2023reward} identify as a source of unfair algorithm comparisons, since weights tuned with one
algorithm tend to favor it.

\paragraph{Reward function.}
All custom rewards are computed per step by a wrapper around the environment (\code{common/rewards.py}) from the
reward components that \env{HighwayEnv}~\cite{highway-env} reports, without the environment's normalization:
\begin{equation}
  r_t = w_{\mathrm{speed}}\, r^{\mathrm{speed}}_t + w_{\mathrm{right}}\, r^{\mathrm{right}}_t
        - w_{\mathrm{col}}\, c_t - w_{\mathrm{lane}}\, \ell_t,
  \label{eq:reward}
\end{equation}
where
\begin{itemize}
  \item $r^{\mathrm{speed}}_t \in [0,1]$ is the high-speed component of \env{HighwayEnv}: the forward speed mapped
  linearly from the default range $[20, 30]$\,m/s to $[0,1]$ and clipped;
  \item $r^{\mathrm{right}}_t \in [0,1]$ is the right-lane component: the index of the target lane divided by the
  number of lanes minus one, so 1 in the rightmost lane and 0 in the leftmost; it encodes the keep-right rule;
  \item $c_t \in \{0,1\}$ indicates a collision, after which the episode terminates;
  \item $\ell_t \in \{0,1\}$ is 1 if the vehicle's actual lane index changed during the step (a completed lane change,
  not the chosen action). This term is applied by the wrapper because \env{HighwayEnv}~1.12.1 ignores the
  \code{lane\_change\_reward} configuration key.
\end{itemize}
We do not use the environment's reward normalization (\cref{eq:default-reward}). It maps the reward to $[0,1]$,
so every step earns a non-negative reward, including the step of a collision. The only deterrent against crashing is
then the survival reward forgone after the crash, which depends on when the crash happens rather than on how large it
is; a larger collision weight only raises this survival bonus and shrinks the speed signal relative to it. Without
normalization, $w_{\mathrm{col}}$ is a direct penalty and the survival bonus disappears.

\paragraph{Weights.}
\Cref{tab:reward-weights} lists the three custom variants and the assignment's labels: Reward~A is efficiency,
Reward~B safety and Reward~C balanced. The weights were set as round values that express the
intended priorities: the speed weight decreases and the collision weight increases from efficiency to safety, and
the keep-right and lane-change terms are absent from the efficiency variant and grow towards safety. In the safety
variant the lane-change penalty (0.5) exceeds the largest reward of a single step ($0.2 + 0.1 = 0.3$), so a safety
agent should change lanes only rarely, even when a lane change would help in dense traffic. The weights are defined
in \code{REWARDS} in \code{common/config.py}.

\begin{table}[H]
  \centering
  \small
  \caption{Reward weights of \cref{eq:reward}. The default configuration is the normalized reward of
  \env{HighwayEnv}, unchanged.}
  \label{tab:reward-weights}
  \begin{tabularx}{\linewidth}{@{}llccccX@{}}
    \toprule
    Variant & Assignment & $w_{\mathrm{speed}}$ & $w_{\mathrm{right}}$ & $w_{\mathrm{col}}$ & $w_{\mathrm{lane}}$ & Intended behavior \\
    \midrule
    Default    & default & \multicolumn{4}{c}{\env{HighwayEnv} default (normalized)} & Reference \\
    Efficiency & Reward A & 0.8 & 0.0 & 15 & 0.0 & Speed first, lane position irrelevant, tolerant of risk \\
    Safety     & Reward B & 0.2 & 0.1 & 30 & 0.5 & Weak speed incentive, harsh collision penalty, costly lane changes \\
    Balanced   & Reward C & 0.4 & 0.1 & 20 & 0.1 & Moderate speed, keep right, mild lane-change penalty \\
    \bottomrule
  \end{tabularx}
\end{table}

\subsection{Sanity check of the weights}\label{sec:reward-check}

\paragraph{Preference ordering and indifference point.}
Knox et al.~\cite{knox2023reward} propose sanity checks for driving rewards based on a few idealized trajectories.
The preference-ordering test requires that a trajectory that crashes halfway is less preferred than a safe one that
does not drive, and the safe one less than a successful trip. The indifference-point test then asks at which success
probability $p$ a gamble between success and crash is worth exactly as much as the safe alternative:
\begin{equation}
  G(\tau_{\mathrm{slow}}) = p\, G(\tau_{\mathrm{succ}}) + (1-p)\, G(\tau_{\mathrm{crash}})
  \quad\Longrightarrow\quad
  p = \frac{G(\tau_{\mathrm{slow}}) - G(\tau_{\mathrm{crash}})}{G(\tau_{\mathrm{succ}}) - G(\tau_{\mathrm{crash}})}.
  \label{eq:indifference}
\end{equation}
A policy that crashes more often than $1-p$ is worse, under the reward, than refusing to drive fast at all, so $p$
measures the risk the reward tolerates. Because the ego vehicle in \env{HighwayEnv} cannot stay parked, we replace
the idle trajectory of Knox et al.\ with slow safe driving. With episodes of $T = 30$ steps (30\,s at one decision per
second), all in the right lane and without lane changes, the idealized trajectories are: success, $T$ steps at
30\,m/s ($r^{\mathrm{speed}} = 1$); slow, $T$ steps at 20\,m/s ($r^{\mathrm{speed}} = 0$); crash, a collision at step
$T/2$ at 25\,m/s ($r^{\mathrm{speed}} = 0.5$). Their undiscounted returns are
\begin{equation}
  G_{\mathrm{succ}} = T\,(w_{\mathrm{speed}} + w_{\mathrm{right}}), \qquad
  G_{\mathrm{slow}} = T\, w_{\mathrm{right}}, \qquad
  G_{\mathrm{crash}} = \tfrac{T}{2}\,(0.5\, w_{\mathrm{speed}} + w_{\mathrm{right}}) - w_{\mathrm{col}}.
  \label{eq:idealized-returns}
\end{equation}

\begin{table}[H]
  \centering
  \small
  \caption{Returns of the idealized trajectories and analytic indifference points ($T = 30$, no discounting).}
  \label{tab:reward-analytic}
  \begin{tabular}{@{}lccccc@{}}
    \toprule
    Variant & $G_{\mathrm{crash}}$ & $G_{\mathrm{slow}}$ & $G_{\mathrm{succ}}$ & Ordering & Analytic $p$ \\
    \midrule
    Default    & 12.3    & 22.0 & 30.0 & holds & 0.55 \\
    Efficiency & $-9.0$  & 0.0 & 24.0 & holds & 0.27 \\
    Balanced   & $-15.5$ & 3.0 & 15.0 & holds & 0.61 \\
    Safety     & $-27.0$ & 3.0 & 9.0  & holds & 0.83 \\
    \bottomrule
  \end{tabular}
\end{table}

All three variants pass the ordering test (\cref{tab:reward-analytic}), and their indifference points increase from
efficiency to safety, so the intended gradient of caution is reflected in the reward and not only in the names. The
default reward passes as well, but its crash costs only the survival reward forgone after it. Analytically, much
smaller collision penalties would already keep the ordering ($w_{\mathrm{col}} > 6$ for efficiency, $> 1.5$ for
balanced, $> 0$ for safety); the binding case is the empirical check below at low density, where crashes happen later
and therefore cost less relative to the reward already collected. Among round values, 15, 20 and 30 are the smallest
that pass this check at every density while keeping the gradient between the variants.

\paragraph{Empirical check on scripted policies.}
The idealized trajectories ignore when crashes actually happen and which lane a slow vehicle uses, so we repeated the
check on simulated episodes of the four scripted policies of \cref{sec:baseline}, 50 episodes per policy and
density, with a 2,000-resample bootstrap; no agent was trained for this check. $G_{\mathrm{crash}}$ pools the crashed
episodes of all policies, $G_{\mathrm{slow}}$ comes from the slower policy, and $G_{\mathrm{succ}}$ from non-crashed
episodes of the faster policy (of the idle policy where fewer than five exist).

\begin{table}[H]
  \centering
  \small
  \caption{Empirical indifference points on scripted policies (undiscounted returns; 5--95\% bootstrap intervals in
  brackets).}
  \label{tab:reward-empirical}
  \begin{tabular}{@{}lcccc@{}}
    \toprule
    Variant & Low density & Medium (training) & High density & Analytic \\
    \midrule
    Efficiency & 0.11 (0.06--0.15) & 0.36 (0.33--0.38) & 0.46 (0.45--0.47) & 0.27 \\
    Balanced   & 0.55 (0.53--0.56) & 0.73 (0.72--0.75) & 0.74 (0.72--0.76) & 0.61 \\
    Safety     & 0.82 (0.81--0.83) & 0.90 (0.89--0.92) & 0.89 (0.88--0.91) & 0.83 \\
    \bottomrule
  \end{tabular}
\end{table}

With undiscounted returns and with $\gam = 0.99$, the ordering holds for every variant at every density, and the
empirical indifference points keep the order efficiency $<$ balanced $<$ safety with non-overlapping bootstrap
intervals (\cref{tab:reward-empirical}). The deviations from the analytic values follow from the crash time (later at
low density, earlier at high density) and from the success reference, which at medium and high density is the idle
policy at 25\,m/s rather than 30\,m/s. With $\gam = 0.99$ the empirical indifference points change by at most 0.04
(efficiency 0.07 / 0.34 / 0.46, balanced 0.53 / 0.73 / 0.75, safety 0.81 / 0.90 / 0.90 at low / medium / high
density). With $\gam = 0.9$ or $0.8$ a late crash is discounted so strongly that the
efficiency and balanced variants fail the ordering at some densities, and only the safety variant passes everywhere.
We therefore train both algorithms with $\gam \geq 0.99$.

\paragraph{What the check does not guarantee.}
The check shows that the rewards rank idealized and scripted outcomes in the intended order; it does not show that a
learned policy will avoid collisions. A collision penalty only changes the expected return, and an agent can still
accept a crash risk below $1-p$, or fail to find a safe policy within the training budget. Returns are also not
comparable across variants, since a collision costs 15, 20 or 30, so the variants are compared through driving
behavior (collision rate, speed, completed episodes, lane use) rather than return.
